\documentclass[10pt,letterpaper]{article}
\usepackage{cornell-notes}
\usepackage{mathtools}

\title{Chapter 9: Root Finding and Nonlinear Sets of Equations}
\author{}
\date{}

\setDocTitle{Chapter 9: Root Finding and Nonlinear Sets of Equations}
\setDocSubtitle{Numerical Methods}
\setDocVersion{v1.0}
\setDocStatus{Study Companion}
\setDocOwner{Numerical Methods Cornell Notes}
\setDocAuthor{}
\setDocDate{\today}
\setDocSummary{Cornell-format study and review notes for Chapter 9: Root Finding and Nonlinear Sets of Equations.}

\setCornellCollection{Numerical Methods Cornell Notes}
\setCornellUnitType{Chapter}
\setCornellUnitNumber{9}
\setCornellUnitTitle{Root Finding and Nonlinear Sets of Equations}
\setCornellCompanionLabel{Study and review companion}

\hypersetup{pdftitle={Chapter 9: Root Finding and Nonlinear Sets of Equations},pdfsubject={Cornell notes},pdfauthor={}}

\begin{document}
\maketitle

\section*{Chapter Overview}
This chapter organizes the mathematical and computational ideas behind root finding and nonlinear sets of equations. The notes preserve the numbered section sequence while emphasizing method selection, explicit assumptions, numerical reliability, cost, and verification.

\section*{Learning Objectives}
Explain the governing problem class; compare Bracketing and Bisection, Secant Method, False Position Method, and Ridders' Method, Van Wijngaarden-Dekker-Brent Method, and Newton-Raphson Method Using Derivative; design defensible stopping and verification checks; distinguish problem conditioning from algorithmic stability.

\section*{Prerequisites}
Continuity, derivatives, linear systems, norms, and iterative convergence.

\section*{Operational Workflow}
Establish existence or a bracket where possible; scale variables and residuals; choose a safeguarded scalar or system method; monitor step and residual; diagnose singular derivatives or Jacobians.

\subsection*{Formula and notation anchor}
\[
x_{k+1}=x_k-\frac{f(x_k)}{f'(x_k)}
\]
Near a simple root and with a good initial estimate, Newton's method is locally quadratic; safeguards are needed away from that region.

\begin{warnbox}{Numerical warning}
Unbracketed divergence, vanishing derivatives, multiple roots, discontinuities mistaken for roots, poor scaling, singular Jacobians, and stopping on a small step with a large residual.
\end{warnbox}

\section{9.0 Introduction}
\begin{cornell}
\noterow{What is the central idea?}{Root finding combines a mathematical existence model with an iteration, a stopping rule, and safeguards against nonconvergence or false convergence.}
\noterow{How should it be used?}{For Introduction, begin from explicit inputs, outputs, preconditions, and representation choices.  Establish existence or a bracket where possible; scale variables and residuals; choose a safeguarded scalar or system method; monitor step and residual; diagnose singular derivatives or Jacobians.}
\noterow{Cost and reliability}{Analyze arithmetic work, storage, data movement, and convergence separately.  Relevant failure pressures include unbracketed divergence, vanishing derivatives, multiple roots, discontinuities mistaken for roots, poor scaling, singular jacobians, and stopping on a small step with a large residual.  Use scaled tolerances and report both success criteria and diagnostic evidence.}
\noterow{Implementation checklist}{\begin{itemize}
\item Validate dimensions, domains, ordering, and structural assumptions.
\item Guard small divisors, invalid states, overflow, underflow, and nonfinite values.
\item Make mutation, workspace, indexing, and termination behavior explicit.
\item Test a simple exact case, a difficult valid case, a boundary case, and an expected failure.
\end{itemize}}
\end{cornell}
\begin{summarybox}{Section summary}
Root finding combines a mathematical existence model with an iteration, a stopping rule, and safeguards against nonconvergence or false convergence.  Select this technique only when its assumptions match the data, and accept a result only after an independent residual, error estimate, invariant, or refinement check.
\end{summarybox}

\section{9.1 Bracketing and Bisection}
\begin{cornell}
\noterow{What is the central idea?}{A sign-changing interval guarantees a root for a continuous function; bisection halves the interval while preserving that invariant.}
\noterow{How should it be used?}{For Bracketing and Bisection, begin from explicit inputs, outputs, preconditions, and representation choices.  Establish existence or a bracket where possible; scale variables and residuals; choose a safeguarded scalar or system method; monitor step and residual; diagnose singular derivatives or Jacobians.}
\noterow{Quantitative anchor}{\[
b_k-a_k=2^{-k}(b_0-a_0)
\]
State the dimensions, scaling, and validity assumptions before using this relation.  Check the resulting residual or invariant rather than trusting an iteration count alone.}
\noterow{Cost and reliability}{Analyze arithmetic work, storage, data movement, and convergence separately.  Relevant failure pressures include unbracketed divergence, vanishing derivatives, multiple roots, discontinuities mistaken for roots, poor scaling, singular jacobians, and stopping on a small step with a large residual.  Use scaled tolerances and report both success criteria and diagnostic evidence.}
\noterow{Implementation checklist}{\begin{itemize}
\item Validate dimensions, domains, ordering, and structural assumptions.
\item Guard small divisors, invalid states, overflow, underflow, and nonfinite values.
\item Make mutation, workspace, indexing, and termination behavior explicit.
\item Test a simple exact case, a difficult valid case, a boundary case, and an expected failure.
\end{itemize}}
\end{cornell}
\begin{summarybox}{Section summary}
A sign-changing interval guarantees a root for a continuous function; bisection halves the interval while preserving that invariant.  Select this technique only when its assumptions match the data, and accept a result only after an independent residual, error estimate, invariant, or refinement check.
\end{summarybox}

\section{9.2 Secant Method, False Position Method, and Ridders' Method}
\begin{cornell}
\noterow{What is the central idea?}{Derivative-free methods interpolate function values, differing in how strongly they preserve a bracket and how rapidly they converge.}
\noterow{How should it be used?}{For Secant Method, False Position Method, and Ridders' Method, begin from explicit inputs, outputs, preconditions, and representation choices.  Establish existence or a bracket where possible; scale variables and residuals; choose a safeguarded scalar or system method; monitor step and residual; diagnose singular derivatives or Jacobians.}
\noterow{Quantitative anchor}{\[
x_{k+1}=x_k-f(x_k)\frac{x_k-x_{k-1}}{f(x_k)-f(x_{k-1})}
\]
State the dimensions, scaling, and validity assumptions before using this relation.  Check the resulting residual or invariant rather than trusting an iteration count alone.}
\noterow{Cost and reliability}{Analyze arithmetic work, storage, data movement, and convergence separately.  Relevant failure pressures include unbracketed divergence, vanishing derivatives, multiple roots, discontinuities mistaken for roots, poor scaling, singular jacobians, and stopping on a small step with a large residual.  Use scaled tolerances and report both success criteria and diagnostic evidence.}
\noterow{Implementation checklist}{\begin{itemize}
\item Validate dimensions, domains, ordering, and structural assumptions.
\item Guard small divisors, invalid states, overflow, underflow, and nonfinite values.
\item Make mutation, workspace, indexing, and termination behavior explicit.
\item Test a simple exact case, a difficult valid case, a boundary case, and an expected failure.
\end{itemize}}
\end{cornell}
\begin{summarybox}{Section summary}
Derivative-free methods interpolate function values, differing in how strongly they preserve a bracket and how rapidly they converge.  Select this technique only when its assumptions match the data, and accept a result only after an independent residual, error estimate, invariant, or refinement check.
\end{summarybox}

\section{9.3 Van Wijngaarden-Dekker-Brent Method}
\begin{cornell}
\noterow{What is the central idea?}{Brent-style iteration combines bisection reliability with secant or inverse-interpolation speed, accepting a fast step only when it is safe.}
\noterow{How should it be used?}{For Van Wijngaarden-Dekker-Brent Method, begin from explicit inputs, outputs, preconditions, and representation choices.  Establish existence or a bracket where possible; scale variables and residuals; choose a safeguarded scalar or system method; monitor step and residual; diagnose singular derivatives or Jacobians.}
\noterow{Quantitative lens}{Identify the governing equation, recurrence, probability law, or geometric predicate for Van Wijngaarden-Dekker-Brent Method.  Define every variable and scale; then derive a residual, error estimate, or invariant that can be checked independently.}
\noterow{Cost and reliability}{Analyze arithmetic work, storage, data movement, and convergence separately.  Relevant failure pressures include unbracketed divergence, vanishing derivatives, multiple roots, discontinuities mistaken for roots, poor scaling, singular jacobians, and stopping on a small step with a large residual.  Use scaled tolerances and report both success criteria and diagnostic evidence.}
\noterow{Implementation checklist}{\begin{itemize}
\item Validate dimensions, domains, ordering, and structural assumptions.
\item Guard small divisors, invalid states, overflow, underflow, and nonfinite values.
\item Make mutation, workspace, indexing, and termination behavior explicit.
\item Test a simple exact case, a difficult valid case, a boundary case, and an expected failure.
\end{itemize}}
\end{cornell}
\begin{summarybox}{Section summary}
Brent-style iteration combines bisection reliability with secant or inverse-interpolation speed, accepting a fast step only when it is safe.  Select this technique only when its assumptions match the data, and accept a result only after an independent residual, error estimate, invariant, or refinement check.
\end{summarybox}

\section{9.4 Newton-Raphson Method Using Derivative}
\begin{cornell}
\noterow{What is the central idea?}{Newton iteration linearizes the function at the current estimate and converges rapidly near a simple root when the derivative is reliable.}
\noterow{How should it be used?}{For Newton-Raphson Method Using Derivative, begin from explicit inputs, outputs, preconditions, and representation choices.  Establish existence or a bracket where possible; scale variables and residuals; choose a safeguarded scalar or system method; monitor step and residual; diagnose singular derivatives or Jacobians.}
\noterow{Quantitative anchor}{\[
x_{k+1}=x_k-\frac{f(x_k)}{f'(x_k)}
\]
State the dimensions, scaling, and validity assumptions before using this relation.  Check the resulting residual or invariant rather than trusting an iteration count alone.}
\noterow{Cost and reliability}{Analyze arithmetic work, storage, data movement, and convergence separately.  Relevant failure pressures include unbracketed divergence, vanishing derivatives, multiple roots, discontinuities mistaken for roots, poor scaling, singular jacobians, and stopping on a small step with a large residual.  Use scaled tolerances and report both success criteria and diagnostic evidence.}
\noterow{Implementation checklist}{\begin{itemize}
\item Validate dimensions, domains, ordering, and structural assumptions.
\item Guard small divisors, invalid states, overflow, underflow, and nonfinite values.
\item Make mutation, workspace, indexing, and termination behavior explicit.
\item Test a simple exact case, a difficult valid case, a boundary case, and an expected failure.
\end{itemize}}
\end{cornell}
\begin{summarybox}{Section summary}
Newton iteration linearizes the function at the current estimate and converges rapidly near a simple root when the derivative is reliable.  Select this technique only when its assumptions match the data, and accept a result only after an independent residual, error estimate, invariant, or refinement check.
\end{summarybox}

\section{9.5 Roots of Polynomials}
\begin{cornell}
\noterow{What is the central idea?}{Polynomial root finding must address scaling, complex conjugate structure, deflation error, multiplicity, and sensitivity of clustered roots.}
\noterow{How should it be used?}{For Roots of Polynomials, begin from explicit inputs, outputs, preconditions, and representation choices.  Establish existence or a bracket where possible; scale variables and residuals; choose a safeguarded scalar or system method; monitor step and residual; diagnose singular derivatives or Jacobians.}
\noterow{Quantitative lens}{Identify the governing equation, recurrence, probability law, or geometric predicate for Roots of Polynomials.  Define every variable and scale; then derive a residual, error estimate, or invariant that can be checked independently.}
\noterow{Cost and reliability}{Analyze arithmetic work, storage, data movement, and convergence separately.  Relevant failure pressures include unbracketed divergence, vanishing derivatives, multiple roots, discontinuities mistaken for roots, poor scaling, singular jacobians, and stopping on a small step with a large residual.  Use scaled tolerances and report both success criteria and diagnostic evidence.}
\noterow{Implementation checklist}{\begin{itemize}
\item Validate dimensions, domains, ordering, and structural assumptions.
\item Guard small divisors, invalid states, overflow, underflow, and nonfinite values.
\item Make mutation, workspace, indexing, and termination behavior explicit.
\item Test a simple exact case, a difficult valid case, a boundary case, and an expected failure.
\end{itemize}}
\end{cornell}
\begin{summarybox}{Section summary}
Polynomial root finding must address scaling, complex conjugate structure, deflation error, multiplicity, and sensitivity of clustered roots.  Select this technique only when its assumptions match the data, and accept a result only after an independent residual, error estimate, invariant, or refinement check.
\end{summarybox}

\section{9.6 Newton-Raphson Method for Nonlinear Systems of Equations}
\begin{cornell}
\noterow{What is the central idea?}{A nonlinear system Newton step solves a Jacobian linear system for a correction rather than dividing by one derivative.}
\noterow{How should it be used?}{For Newton-Raphson Method for Nonlinear Systems of Equations, begin from explicit inputs, outputs, preconditions, and representation choices.  Establish existence or a bracket where possible; scale variables and residuals; choose a safeguarded scalar or system method; monitor step and residual; diagnose singular derivatives or Jacobians.}
\noterow{Quantitative anchor}{\[
J(x_k)\,\Delta x_k=-F(x_k),\qquad x_{k+1}=x_k+\Delta x_k
\]
State the dimensions, scaling, and validity assumptions before using this relation.  Check the resulting residual or invariant rather than trusting an iteration count alone.}
\noterow{Cost and reliability}{Analyze arithmetic work, storage, data movement, and convergence separately.  Relevant failure pressures include unbracketed divergence, vanishing derivatives, multiple roots, discontinuities mistaken for roots, poor scaling, singular jacobians, and stopping on a small step with a large residual.  Use scaled tolerances and report both success criteria and diagnostic evidence.}
\noterow{Implementation checklist}{\begin{itemize}
\item Validate dimensions, domains, ordering, and structural assumptions.
\item Guard small divisors, invalid states, overflow, underflow, and nonfinite values.
\item Make mutation, workspace, indexing, and termination behavior explicit.
\item Test a simple exact case, a difficult valid case, a boundary case, and an expected failure.
\end{itemize}}
\end{cornell}
\begin{summarybox}{Section summary}
A nonlinear system Newton step solves a Jacobian linear system for a correction rather than dividing by one derivative.  Select this technique only when its assumptions match the data, and accept a result only after an independent residual, error estimate, invariant, or refinement check.
\end{summarybox}

\section{9.7 Globally Convergent Methods for Nonlinear Systems of Equations}
\begin{cornell}
\noterow{What is the central idea?}{Line searches or trust-like safeguards reduce a merit function when the full Newton step is not globally reliable.}
\noterow{How should it be used?}{For Globally Convergent Methods for Nonlinear Systems of Equations, begin from explicit inputs, outputs, preconditions, and representation choices.  Establish existence or a bracket where possible; scale variables and residuals; choose a safeguarded scalar or system method; monitor step and residual; diagnose singular derivatives or Jacobians.}
\noterow{Quantitative anchor}{\[
J(x_k)\,\Delta x_k=-F(x_k),\qquad x_{k+1}=x_k+\Delta x_k
\]
State the dimensions, scaling, and validity assumptions before using this relation.  Check the resulting residual or invariant rather than trusting an iteration count alone.}
\noterow{Cost and reliability}{Analyze arithmetic work, storage, data movement, and convergence separately.  Relevant failure pressures include unbracketed divergence, vanishing derivatives, multiple roots, discontinuities mistaken for roots, poor scaling, singular jacobians, and stopping on a small step with a large residual.  Use scaled tolerances and report both success criteria and diagnostic evidence.}
\noterow{Implementation checklist}{\begin{itemize}
\item Validate dimensions, domains, ordering, and structural assumptions.
\item Guard small divisors, invalid states, overflow, underflow, and nonfinite values.
\item Make mutation, workspace, indexing, and termination behavior explicit.
\item Test a simple exact case, a difficult valid case, a boundary case, and an expected failure.
\end{itemize}}
\end{cornell}
\begin{summarybox}{Section summary}
Line searches or trust-like safeguards reduce a merit function when the full Newton step is not globally reliable.  Select this technique only when its assumptions match the data, and accept a result only after an independent residual, error estimate, invariant, or refinement check.
\end{summarybox}

\section*{Method comparison}
\begin{center}
\small
\begin{tabularx}{\linewidth}{>{\bfseries\raggedright\arraybackslash}p{0.22\linewidth}XX}
\toprule
Method or lens & Best use & Main caution \\
\midrule
Bisection & Continuous sign-changing bracket & Guaranteed linear convergence \\
Brent-style method & Bracketed scalar root & Reliability plus superlinear steps \\
Newton & Reliable derivatives and nearby start & Very fast locally; can diverge \\
\bottomrule
\end{tabularx}
\end{center}

\section*{Worked example}
\begin{exambox}{Independent worked example}
For $f(x)=x^2-2$ and $x_0=1.5$, Newton gives $x_1=1.5-(0.25/3)=1.4166667$ and $x_2\approx1.4142157$.  The residual $|x_2^2-2|$ is about $6\times10^{-6}$.  Stop only when both residual and step meet scaled tolerances.
\end{exambox}

\section*{Chapter synthesis}
\begin{summarybox}{Chapter synthesis}
The chapter's methods should be viewed as a decision system rather than a list of routines.  Begin with the mathematical problem and its conditioning, exploit structure, select an algorithm with compatible stability and cost, and verify the computed answer with evidence that is independent of the internal iteration.  The recurring implementation discipline is: establish existence or a bracket where possible; scale variables and residuals; choose a safeguarded scalar or system method; monitor step and residual; diagnose singular derivatives or jacobians.
\end{summarybox}

\subsection*{Common mistakes and numerical hazards}
\begin{warnbox}{Common mistakes and numerical hazards}
Unbracketed divergence, vanishing derivatives, multiple roots, discontinuities mistaken for roots, poor scaling, singular Jacobians, and stopping on a small step with a large residual.  A second recurring mistake is reporting only a computed value: always retain tolerances, iteration counts, residuals, refinement behavior, and relevant structural checks.
\end{warnbox}

\subsection*{Retrieval practice}
\textbf{Questions}
\begin{enumerate}
\item What problem does Bracketing and Bisection address, and which assumption is easiest to violate?
\item What problem does Secant Method, False Position Method, and Ridders' Method address, and which assumption is easiest to violate?
\item What problem does Van Wijngaarden-Dekker-Brent Method address, and which assumption is easiest to violate?
\item What problem does Newton-Raphson Method for Nonlinear Systems of Equations address, and which assumption is easiest to violate?
\item What problem does Globally Convergent Methods for Nonlinear Systems of Equations address, and which assumption is easiest to violate?
\item How do conditioning, stability, and the final residual play different roles in this chapter?
\end{enumerate}

\textbf{Brief answers}
\begin{enumerate}
\item A sign-changing interval guarantees a root for a continuous function; bisection halves the interval while preserving that invariant. Recheck the section's domain, structure, scaling, and failure diagnostics.
\item Derivative-free methods interpolate function values, differing in how strongly they preserve a bracket and how rapidly they converge. Recheck the section's domain, structure, scaling, and failure diagnostics.
\item Brent-style iteration combines bisection reliability with secant or inverse-interpolation speed, accepting a fast step only when it is safe. Recheck the section's domain, structure, scaling, and failure diagnostics.
\item A nonlinear system Newton step solves a Jacobian linear system for a correction rather than dividing by one derivative. Recheck the section's domain, structure, scaling, and failure diagnostics.
\item Line searches or trust-like safeguards reduce a merit function when the full Newton step is not globally reliable. Recheck the section's domain, structure, scaling, and failure diagnostics.
\item Conditioning describes sensitivity of the mathematical problem, stability describes error propagation by the algorithm, and the residual measures satisfaction of the computed equations; none is a universal substitute for the others.
\end{enumerate}

\subsection*{Mastery checklist}
\begin{itemize}
\item[$\square$] I can explain and select Bracketing and Bisection without relying on a memorized code listing.
\item[$\square$] I can explain and select Secant Method, False Position Method, and Ridders' Method without relying on a memorized code listing.
\item[$\square$] I can explain and select Van Wijngaarden-Dekker-Brent Method without relying on a memorized code listing.
\item[$\square$] I can explain and select Newton-Raphson Method Using Derivative without relying on a memorized code listing.
\item[$\square$] I can explain and select Roots of Polynomials without relying on a memorized code listing.
\item[$\square$] I can state a scaled stopping rule and an independent verification check.
\item[$\square$] I can identify at least one conditioning risk and one algorithmic stability risk.
\item[$\square$] I can estimate time, storage, and data-movement costs at the appropriate level.
\end{itemize}

\end{document}

